%
% anomalien.tex -- Die Kepler-Gleichung, Ellipse mit Hilfskreis,
%                  exzentrische Anomalie und überstrichene Fläche
%
% (c) 2026 Gian Cavegn, OST Ostschweizer Fachhochschule
%
\documentclass[tikz]{standalone}
\usepackage{amsmath}
\usepackage{times}
\usepackage{txfonts}
\usetikzlibrary{arrows,intersections,math,calc}
\begin{document}
\def\skala{1.15}
\begin{tikzpicture}[>=latex,scale=\skala]
% a = 3, epsilon = 0.6, b = 2.4, exzentrische Anomalie E = 50 Grad
% ueberstrichene Flaeche
\fill[gray!20] (1.8,0) -- (3,0) arc[start angle=0,end angle=50,x radius=3,y radius=2.4] -- cycle;
% Hilfskreis und Ellipse
\draw[dashed] (0,0) circle (3);
\draw[thick] (0,0) ellipse (3 and 2.4);
\draw (-3.4,0) -- (3.4,0);
% Projektion und Fahrstrahlen
\draw[dashed] (1.928,0) -- (1.928,2.298);
\draw (0,0) -- (1.928,2.298);
\draw (1.8,0) -- (1.928,1.838);
% Winkel E
\draw (0.9,0) arc[start angle=0,end angle=50,radius=0.9];
\node at (1.15,0.38) {$E$};
% Punkte
\fill (0,0) circle (1.5pt);
\node at (0,0) [below left] {$O$};
\fill (1.8,0) circle (2pt);
\node at (1.8,0) [below right] {$S$};
\fill (1.928,1.838) circle (1.5pt);
\node at (1.928,1.838) [right] {$P$};
\fill (1.928,2.298) circle (1.5pt);
\node at (1.928,2.298) [above right] {$Q$};
\node at (1.5,1.55) {$a$};
\end{tikzpicture}
\end{document}
